% ---------------------------------------------------------------------------
% Chapter 3: Whose Child Is It? Matching and Merging
% A short stop, by design (prompt/prompt_10_MatchingMerging.md, 2026-10-07):
% our event u ubar -> b bbar is a single leading-order 2->2 process with no
% extra-parton subprocess, so no matching or merging is applied to it. The
% chapter explains in two pages what the procedures are, where CMS uses
% them, and why the ledger does not change here. There is no film and no
% standalone script at this stop; the sketches live in chapter/tikz/mm_*.tex.
% ---------------------------------------------------------------------------
\chapter{Whose Child Is It? Matching and Merging}
\chaptermark{Whose Child Is It?}
\label{ch:whose_child}

\begin{journeybox}
When the matrix element and the shower both claim the same child, the family
would be counted twice. Matching and merging decide, emission by emission,
who is the rightful parent. Our mother was born alone, from a single
$2\to2$ scattering, so at this stop nobody disputes her children: it is a
short stop, a look at the map of a road we do not have to take.
\end{journeybox}

\physicsline{double counting between a $2\to3$ matrix element and the first
shower emission, the merging scale, MLM jet matching and the
\code{xqcut}/\code{qcut} parameters, CKKW-L Sudakov reweighting, NLO matching
in MC@NLO and POWHEG, FxFx merging, where CMS uses each and why our event
uses none}

\section{A short stop: why our event needs no matching}
\label{sec:whose_child:why_short}

\Cref{ch:mother_gives_birth} ended with a confession. The standalone shower,
and \PYTHIA inside \CMSSW, produce every parton beyond the $b$ and $\bar b$
from the splitting functions of \cref{eq:mother_gives_birth:ptmax} downward:
the hardest branching of our event, an initial-state gluon at
$\pt=53.9\GeV$, and the mother's own first emission at $34.2$ and $32.1\GeV$
for the twins, were all generated by a formula that is exact only when the
emission is soft or collinear. A wide-angle gluon of $50\GeV$ is neither,
and a generator that cares about it computes it from the $2\to3$ matrix
element $u\bar u\to\bbbar g$ instead. The moment it does so, a question
arises that this chapter is named after: the $2\to2$ matrix element followed
by one shower emission and the $2\to3$ matrix element describe the \emph{same}
final state, and adding the two counts the gluon twice.

We do not have to answer it, for a reason that is worth stating plainly
rather than hiding. Our event is one leading-order process,
\code{HardQCD:qqbar2bbbar}, with no additional subprocess
($u\bar u\to\bbbar g$, $ug\to\bbbar u$, \dots) anywhere in the sample, neither
in the standalone script, nor in the \CMSSW configuration of
\cref{sec:ancestral_home:validation}, nor in the \MG run of
\cref{sec:ancestral_home:madgraph}, which produced $\bbbar$ and nothing
else. There is exactly one parent for every extra parton, the shower, and
nothing to match it to. The price is the one named in
\cref{sec:mother_gives_birth:caveats}: wide-angle emissions above the
starting scale of $113.32\GeV$ are forbidden and those below it are only
logarithmically accurate. For following one $b$ jet of $100\GeV$ through
the detector that price is small, a few percent of the family's momentum
against a seed-to-seed spread of $29\GeV$, and the official CMS QCD
multijet samples made with \PYTHIA alone pay it too. It is \emph{not} small
when the extra jets are the measurement, $\ttbar$+jets, $W/Z$+jets or the
multijet background of a search, and that is where the recipes below are
used. This stop therefore has no film and no script; it has two pictures,
the vocabulary, and the references to go further.

\section{Who owns the extra gluon?}
\label{sec:whose_child:who_owns}

Guess first: if we had included $u\bar u\to\bbbar g$ in \cref{ch:ancestral_home},
how many of the $38$ branchings of our event would it have produced?
\Cref{fig:whose_child:double_counting} gives the answer with the merging
scale CMS uses most often.

\begin{figure}[htbp]\centering
\resizebox{\textwidth}{!}{% ---------------------------------------------------------------------------
% Matching and merging in one picture. (a) The same final state, b bbar + one
% gluon, reached twice: by the 2->2 matrix element followed by one shower
% emission, and by the 2->3 matrix element. (b) The transverse-momentum axis
% of the extra parton with the two descriptions, the merging scale, and the
% branchings of OUR event (winners of the -v log of
% standalone_parton_showering.py, seed 1), none of which was ever at risk of
% being counted twice because the record has no 2->3 matrix element.
% ---------------------------------------------------------------------------
\begin{tikzpicture}[x=1cm,y=1cm,
  every node/.style={font=\footnotesize},
  fermion/.style={thick,postaction={decorate},decoration={markings,mark=at position 0.55 with {\arrow[scale=1.1]{Stealth}}}},
  antifermion/.style={thick,postaction={decorate},decoration={markings,mark=at position 0.45 with {\arrowreversed[scale=1.1]{Stealth}}}},
  gluon/.style={thick,decorate,decoration={coil,aspect=0,segment length=4pt,amplitude=1.8pt}},
  vertex/.style={fill=black,circle,inner sep=1.5pt},
  lab/.style={font=\scriptsize,align=center},
  head/.style={font=\small\bfseries},
  blob/.style={draw,fill=gray!25,circle,minimum size=7mm,inner sep=0pt},
  me/.style={fill=blue!10,draw=blue!60!black,rounded corners=2pt,inner sep=3pt,font=\scriptsize,align=center},
  ps/.style={fill=orange!15,draw=orange!70!black,rounded corners=2pt,inner sep=3pt,font=\scriptsize,align=center},
  tick/.style={very thick}]
% ---------------- (a) the same child, claimed twice ------------------------------------------------
\node[head,anchor=west] at (-0.3,3.3) {(a) one child, two parents};
% left: 2->2 ME + shower emission
\begin{scope}
\coordinate (V1) at (1.2,0.9); \coordinate (V2) at (2.9,0.9);
\draw[fermion] (0,2.0) -- (V1) node[pos=0,above left,lab] {$u$};
\draw[antifermion] (0,-0.2) -- (V1) node[pos=0,below left,lab] {$\bar u$};
\draw[gluon] (V1) -- (V2);
\coordinate (E) at (3.9,1.7);
\draw[fermion] (V2) -- (E); \draw[fermion] (E) -- (4.9,2.3) node[right,lab] {$b$};
\draw[gluon,orange!80!black] (E) -- (5.0,1.2) node[right,lab,text=orange!80!black] {$g$ (shower)};
\draw[antifermion] (V2) -- (4.9,-0.2) node[right,lab] {$\bar b$};
\node[vertex] at (V1) {}; \node[vertex] at (V2) {}; \node[vertex,fill=orange!80!black] at (E) {};
\node[me,anchor=north west] at (-0.2,-0.9) {$|\mathcal M_{2\to2}|^2$ \textbf{+} shower kernel\\$\dfrac{\alphas}{2\pi}\dfrac{\mathrm d\pt^2}{\pt^2}\,P_{qq}(z)$};
\node[lab,text=black!60] at (2.4,-2.6) {chapter 2: $113.32\GeV$ bounds the gluon};
\end{scope}
\node[font=\large] at (6.9,0.9) {$\neq$};
% right: 2->3 ME
\begin{scope}[xshift=8.2cm]
\coordinate (V1) at (1.2,0.9); \coordinate (V2) at (2.9,0.9);
\draw[fermion] (0,2.0) -- (V1) node[pos=0,above left,lab] {$u$};
\draw[antifermion] (0,-0.2) -- (V1) node[pos=0,below left,lab] {$\bar u$};
\draw[gluon] (V1) -- (V2);
\coordinate (E) at (3.9,1.7);
\draw[fermion] (V2) -- (E); \draw[fermion] (E) -- (4.9,2.3) node[right,lab] {$b$};
\draw[gluon,blue!60!black] (E) -- (5.0,1.2) node[right,lab,text=blue!60!black] {$g$ (matrix element)};
\draw[antifermion] (V2) -- (4.9,-0.2) node[right,lab] {$\bar b$};
\node[vertex] at (V1) {}; \node[vertex] at (V2) {}; \node[vertex,fill=blue!60!black] at (E) {};
\node[lab,text=black!60,anchor=north west] at (-0.2,-1.0) {+ the gluon from the $\bar b$, from the\\initial state, and all interferences};
\node[me,anchor=north west] at (-0.2,-1.9) {$|\mathcal M_{2\to3}|^2$: exact at tree level,\\ singular as $\pt\to0$, no Sudakov};
\end{scope}
% ---------------- (b) the pT axis of the extra parton ---------------------------------------------
\begin{scope}[yshift=-7.2cm]
\node[head,anchor=west] at (-0.3,3.5) {(b) the merging scale divides the \pt of the extra parton, and where our event's branchings fall};
% axis: log-ish, 0.5 .. 130 GeV mapped to 0..14 cm by hand
\draw[-{Stealth},thick] (0,0) -- (14.6,0) node[right,lab] {$\pt$ of the extra parton};
\foreach \x/\l in {0.4/1,3.3/3,5.4/10,8.1/30,10.9/100} {\draw (\x,0.1) -- (\x,-0.1); \node[lab,below] at (\x,-0.1) {$\l\GeV$};}
% regions
\fill[orange!15] (0,1.5) rectangle (8.1,2.5);
\fill[blue!10] (8.1,1.5) rectangle (11.5,2.5);
\node[ps,draw=none,fill=none] at (4.0,2.0) {\textbf{shower territory}: $\pt<Q_{\mathrm{cut}}$\\Sudakov-resummed, many soft and collinear partons};
\node[me,draw=none,fill=none] at (9.8,2.0) {\textbf{matrix-element territory}\\$\pt>Q_{\mathrm{cut}}$: exact, wide angle};
\draw[dashed,thick,red!70!black] (8.1,-0.3) -- (8.1,2.7);
\node[lab,text=red!70!black,anchor=east] at (8.0,2.9) {merging scale $Q_{\mathrm{cut}}\simeq30\GeV$ (CMS FxFx)};
\draw[dashed,thick,black!60] (11.5,-0.3) -- (11.5,2.7);
\node[lab,text=black!60,anchor=west] at (11.6,2.9) {$\pt^{\max}=113.3\GeV$: chapter 2's wall};
% our event's branchings: ISR 53.9 51.4 42.2 (orange, above axis), FSR 34.2 32.1 25.6 22.1 13.4 10.6 (red), then many below 6
\foreach \x in {9.5,9.4,8.9} {\draw[tick,green!50!black] (\x,0.0) -- (\x,0.6);}
\foreach \x in {8.4,8.2,7.7,7.3,6.1,5.6} {\draw[tick,red!70!black] (\x,0.0) -- (\x,0.6);}
\foreach \x in {4.1,4.0,3.9,3.4,2.5,2.4,2.2,1.9,1.8,1.5,1.2,1.0,0.9,0.7,0.6,0.5,0.4} {\draw[tick,gray] (\x,0.0) -- (\x,0.4);}
\node[lab,text=green!50!black,anchor=south west] at (8.8,0.65) {ISR 53.9, 51.4, 42.2};
\node[lab,text=red!70!black,anchor=south] at (6.5,0.65) {FSR 34.2 ($\bar b$), 32.1 ($b$),\\25.6, 22.1, 13.4, 10.6};
\node[lab,text=gray,anchor=south] at (2.3,0.45) {32 more, all below $6\GeV$};
\node[lab,text=black!65,text width=14.4cm,anchor=north west] at (-0.3,-0.7) {Ticks: the 38 winning branchings of our event, from the \code{-v} log of chapter 2. In a merged sample the five hardest would have been produced by a $2\to3$ or $2\to4$ matrix element and the shower would have been \emph{forbidden} to emit above $Q_{\mathrm{cut}}$; here the shower produced all of them, which is exact only in the collinear and soft limits. For a single $b$ jet of $\pt\approx100\GeV$ the difference is a few percent of the family's momentum, well inside the seed-to-seed spread of $29\GeV$.};
\end{scope}
\end{tikzpicture}
}
\caption[The same gluon from two parents, and the merging scale against the branchings of our event]{(a) The same $\bbbar g$ final state reached by the $2\to2$ matrix element followed by one shower emission, and by the $2\to3$ matrix element; summing them counts the gluon twice. (b) The transverse momentum of the extra parton with the two territories separated by the merging scale, and the $38$ winning branchings of our event from the \code{-v} log of \cref{ch:mother_gives_birth}: five lie above a merging scale of $30\GeV$ and would, in a merged sample, have come from a matrix element.}
\label{fig:whose_child:double_counting}
\end{figure}

The arithmetic of the problem fits in one line. With $\mathcal P(\pt)$ the
shower's probability of an emission at \pt from the Born configuration,
$\Delta$ its Sudakov factor and $\sigma_1$ the $2\to3$ cross section with the
extra parton above \pt, the two descriptions of the emission rate are
\begin{equation}
  \frac{\mathrm d\sigma}{\mathrm d\pt}\bigg|_{\mathrm{shower}}
   = \sigma_0\,\Delta(\pt^{\max},\pt)\,\mathcal P(\pt),
  \qquad
  \frac{\mathrm d\sigma}{\mathrm d\pt}\bigg|_{\mathrm{ME}}
   = \frac{\mathrm d\sigma_1}{\mathrm d\pt},
  \qquad
  \mathcal P(\pt)\xrightarrow{\ \pt\to0\ }\frac{1}{\sigma_0}\frac{\mathrm d\sigma_1}{\mathrm d\pt}.
  \label{eq:whose_child:two_rates}
\end{equation}
The last limit is the whole story: at small \pt the matrix element and the
shower kernel agree, so the shower's Sudakov factor is the right thing to
keep there; at large \pt they differ, the matrix element is right and the
shower has no business emitting. Every recipe of
\cref{sec:whose_child:recipes} is a way of drawing the line between the
two, at a \emph{merging scale} $Q_{\mathrm{cut}}$, so that the sum is
neither double counted nor discontinuous. For our event the shower emitted
on both sides of the line: three initial-state branchings at $53.9$, $51.4$
and $42.2\GeV$ and the twins' first gluons at $34.2$ and $32.1\GeV$ would
have been the matrix element's children in a merged sample with
$Q_{\mathrm{cut}}=30\GeV$, the remaining $33$, all below $26\GeV$ and
$32$ of them below $6\GeV$, the shower's. \PYTHIA's own answer for a plain
$2\to2$ QCD process is cruder than any merging: it either stops the shower
at the factorisation scale, the ``wall'' of \cref{eq:mother_gives_birth:ptmax}
(\code{PTmaxMatch = 1}, the \CMSSW default), or lets it run to the kinematic
limit with a damping factor that mimics the matrix element on average, the
``power shower'' of \cite{Plehn:2005cq,Corke:2010zj}. The exact
matrix-element corrections \PYTHIA does apply, following
\cite{Miu:1998ju,Norrbin:2000uu}, cover colour-singlet production and
resonance decays, not a coloured $2\to2$ final state like ours.

\section{The recipes, one paragraph each}
\label{sec:whose_child:recipes}

\begin{figure}[htbp]\centering
\resizebox{\textwidth}{!}{% ---------------------------------------------------------------------------
% The three families of recipes as flow charts: (a) tree-level merging of jet
% multiplicities (MLM and CKKW-L), (b) NLO matching of one multiplicity
% (MC@NLO and POWHEG), (c) FxFx = (a) applied to (b). Nothing in this figure
% is used for our event; it is the map of what the official samples do.
% ---------------------------------------------------------------------------
\begin{tikzpicture}[x=1cm,y=1cm,
  every node/.style={font=\footnotesize},
  box/.style={draw,rounded corners=2pt,inner sep=4pt,align=center,font=\scriptsize,minimum width=2.6cm},
  me/.style={box,fill=blue!10,draw=blue!60!black},
  ps/.style={box,fill=orange!15,draw=orange!70!black},
  dec/.style={box,fill=red!8,draw=red!60!black},
  result/.style={box,fill=green!10,draw=green!50!black},
  flow/.style={-{Stealth[length=2mm]},thick},
  head/.style={font=\small\bfseries},
  lab/.style={font=\scriptsize,align=center}]
% ---------------- (a) tree-level merging ---------------------------------------------------------
\node[head,anchor=west] at (0,5.4) {(a) tree-level merging (MLM, CKKW-L): several multiplicities, one sample};
\node[me] (me0) at (1.4,4.2) {$u\bar u\to\bbbar$\\$+0$ partons};
\node[me] (me1) at (4.7,4.2) {$u\bar u\to\bbbar g$\\$+1$ parton, $\kt>\code{xqcut}$};
\node[me] (me2) at (8.4,4.2) {$u\bar u\to\bbbar gg,\ \bbbar q\bar q$\\$+2$ partons};
\node[ps] (sh) at (4.6,2.7) {shower each event with the emission\\ scale bounded by the parton it already has};
\node[dec] (mlm) at (3.5,1.3) {MLM: cluster the showered partons into \kt jets at \code{qcut};\\ keep the event only if every ME parton has its own jet\\ (and the highest multiplicity may have extra jets)};
\node[dec] (ckkw) at (11.2,1.3) {CKKW-L: reweight by the Sudakov factors\\ of the shower history the ME would have had,\\ veto shower emissions above $Q_{\mathrm{cut}}$};
\node[result] (out) at (7.6,-0.3) {one sample, exact wide-angle jets above $Q_{\mathrm{cut}}$,\\ resummed below, smooth across the scale};
\draw[flow] (me0) -- (sh); \draw[flow] (me1) -- (sh); \draw[flow] (me2) -- (sh);
\draw[flow] (sh) -- (mlm); \draw[flow] (sh) -- (ckkw);
\draw[flow] (mlm) -- (out); \draw[flow] (ckkw) -- (out);
\node[lab,anchor=west,text=black!60] at (10.5,4.2) {$\sigma_n\propto\alphas^{\,n}$: each extra\\ leg costs a power of $\alphas$,\\ cut off by \code{xqcut} to stay finite};
% ---------------- (b) NLO matching ---------------------------------------------------------------
\begin{scope}[yshift=-7.2cm]
\node[head,anchor=west] at (0,5.4) {(b) NLO matching (MC@NLO, POWHEG)};
\node[me] (born) at (1.6,4.2) {Born $u\bar u\to\bbbar$\\ $+$ one-loop virtual};
\node[me] (real) at (5.1,4.2) {real emission\\ $u\bar u\to\bbbar g$, $ug\to\bbbar u$, \dots};
\node[dec] (sub) at (3.1,2.7) {subtract from the real emission what the shower's\\ first emission would generate anyway (MC@NLO),\\ or generate the hardest emission with the exact\\ real matrix element and let the shower fill in below (POWHEG)};
\node[result] (outb) at (3.1,1.0) {NLO cross section and NLO shape of the first emission;\\ MC@NLO events may carry negative weights, POWHEG's do not};
\draw[flow] (born) -- (sub); \draw[flow] (real) -- (sub); \draw[flow] (sub) -- (outb);
% ---------------- (c) FxFx ---------------------------------------------------------------
\node[head,anchor=west] at (8.6,5.4) {(c) FxFx = (a) on top of (b)};
\node[me] (f0) at (9.6,4.2) {$\bbbar+0$ at NLO};
\node[me] (f1) at (12.6,4.2) {$\bbbar+1$ at NLO};
\node[dec] (fx) at (11.1,2.7) {MC@NLO each, then MLM-style\\ jet matching at $Q_{\mathrm{cut}}$ with\\ Sudakov reweighting between them};
\node[result] (outc) at (11.1,1.0) {NLO accuracy for the 0- and 1-jet rates\\ in one sample (CMS: $\ttbar$, $W/Z$+jets)};
\draw[flow] (f0) -- (fx); \draw[flow] (f1) -- (fx); \draw[flow] (fx) -- (outc);
\end{scope}
\end{tikzpicture}
}
\caption[Tree-level merging, NLO matching and FxFx as flow charts]{The three families of recipes. (a) Tree-level merging combines samples of different parton multiplicities, each cut off by \code{xqcut}, and decides after the shower whether an event is kept (MLM) or how it is weighted (CKKW-L). (b) NLO matching corrects one multiplicity by one order, either by subtracting the shower's first emission from the real-emission matrix element (MC@NLO) or by generating the hardest emission exactly and showering below it (POWHEG). (c) FxFx applies (a) to samples made with (b). None of them touches our event.}
\label{fig:whose_child:recipes}
\end{figure}

\paragraph{Tree-level merging: MLM.} Generate $\bbbar+0$, $+1$, $+2$ partons
at tree level, each with the extra partons separated by more than
\code{xqcut} in \kt so that the $2\to3$ cross section is finite; shower
each event; cluster the showered partons into \kt jets at a scale
\code{qcut} $>$ \code{xqcut}; keep the event only if every matrix-element
parton is matched to its own jet, and for the highest multiplicity also if
there are extra, softer jets. Events in which the shower produced a jet
above \code{qcut} that the matrix element did not are rejected, which is how
the double counting is removed and how the Sudakov suppression of the
matrix-element legs is applied, event by event rather than as a
weight~\cite{Mangano:2006rw,Alwall:2007fs}. In \PYTHIA the step is
\code{JetMatching:scheme = 1} with \code{JetMatching:qCut} and
\code{JetMatching:nJetMax}; CMS uses it for its \MG multijet, Drell--Yan and
$W$+jets leading-order samples with \code{qcut} between $19$ and $30\GeV$.

\paragraph{Tree-level merging: CKKW and CKKW-L.} Instead of rejecting events,
reconstruct for each matrix-element configuration the shower history that
would have produced it, weight the event by the Sudakov factors and
coupling ratios of that history, and veto shower emissions above the
merging scale~\cite{Catani:2001cc}. L\"onnblad's version uses the
shower's own Sudakov factors, obtained by trial showers, so that the
weights are exactly what the shower would have given~\cite{Lonnblad:2001iq,Lonnblad:2011xx};
the unitarised form UMEPS makes the merged inclusive cross section equal
to the lowest-multiplicity one~\cite{Lonnblad:2012ng}, and Sherpa's
implementation with truncated showers is~\cite{Hoeche:2009rj}. All these
are sub-percent-level different for well-separated jets and differ mainly
in how smooth the result is across $Q_{\mathrm{cut}}$.

\paragraph{NLO matching: MC@NLO and POWHEG.} Both raise one multiplicity to
next-to-leading order, which fixes the normalisation to $\sim10\%$
(\cref{sec:ancestral_home:scales}) and gives the hardest emission
its exact shape. MC@NLO subtracts from the real-emission matrix element the
term the shower's first emission generates, so that the sum is NLO without
double counting; the subtraction can be larger than the term, and a fraction
of the events, typically $10$--$30\%$, carries negative
weight~\cite{Frixione:2002ik,Alwall:2014hca}. POWHEG instead generates the
hardest emission with the exact real matrix element and a Sudakov built from
it, then hands the event to a shower restricted below that emission's
scale; the weights are positive but the shower must respect the
\code{SCALUP} bound, which for \PYTHIA's \pt-ordered shower is only
approximately the same variable~\cite{Nason:2004rx,Frixione:2007vw}. CMS's
inclusive $\ttbar$ and single-top samples are POWHEG.

\paragraph{FxFx and the NLO mergings.} FxFx is MLM applied to MC@NLO
samples of different multiplicities, with the Sudakov reweighting between
them done consistently at NLO~\cite{Frederix:2012ps,Frederix:2015eii}; in
\PYTHIA it is \code{JetMatching:doFxFx = on} with \code{JetMatching:qCut}
and \code{JetMatching:nQmatch}, and CMS uses it with $Q_{\mathrm{cut}}=30\GeV$
for $\ttbar$+jets and $W/Z$+jets, where the measurement of
\cite{CMS:2022ilp} compares it with the MLM and other predictions jet by
jet. MEPS@NLO~\cite{Hoeche:2012yf} and UNLOPS~\cite{Lonnblad:2012ix} are the
CKKW-type counterparts; MiNLO and MiNNLO$_{\mathrm{PS}}$
\cite{Hamilton:2012np,Monni:2019whf} reach the same goal from the POWHEG
side and, for colour-singlet production, at NNLO. The review
\cite{Buckley:2011ms} and the \PYTHIA guide \cite{Bierlich:2022pfr},
chapter on matching and merging, are the places to read all of this at
length.

\section{Our jet at this stage: nothing changes, by construction}
\label{sec:whose_child:our_jet}

The ledger receives a row that repeats the previous one. It is there so
that the chain of stages is complete, and as a reminder that for a merged
sample this row would be the first in which the matrix element, not the
shower, decided whether the mother's $34\GeV$ gluon exists.

\begin{ledger}{Whose Child Is It?}
  \ledgerrow{$b$ family after the shower ($\dR<0.4$)}{100.45}{13.52}{3}{--}
  \ledgerrow{$b$ family after matching and merging (none applied)}{100.45}{13.52}{3}{$0$}
\end{ledger}

\section{Further reading}
\label{sec:whose_child:further_reading}
MLM matching is defined in \cite{Mangano:2006rw} and compared with CKKW,
CKKW-L and others in \cite{Alwall:2007fs}; CKKW is \cite{Catani:2001cc},
its dipole-shower form \cite{Lonnblad:2001iq,Lonnblad:2011xx}; MC@NLO is
\cite{Frixione:2002ik}, POWHEG \cite{Nason:2004rx,Frixione:2007vw}, FxFx
\cite{Frederix:2012ps}. The \MG paper \cite{Alwall:2014hca} documents the
settings CMS writes into its gridpacks, and \cite{Buckley:2011ms} is the
general review.

\begin{pitfalls}
\begin{itemize}[nosep]
  \item \textbf{No merging is not an error for our sample.} A \PYTHIA
    \code{HardQCD} sample is a leading-order $2\to2$ sample showered with
    \code{PTmaxMatch = 1}; nothing is missing from it except the physics
    it never claimed, hard wide-angle jets. Do not add \code{JetMatching}
    settings to a sample that has no matrix-element partons to match.
  \item \textbf{\code{qcut} must exceed \code{xqcut}, and both are physics.}
    Changing \code{qcut} by a factor two changes the one-jet rate of an MLM
    sample by several percent; the variation is a systematic uncertainty
    in CMS analyses, not a tuning knob.
  \item \textbf{Negative weights are not a bug.} An MC@NLO or FxFx sample
    with $20\%$ negative weights needs $(1-2\times0.2)^{-2}\approx3$ times
    the events for the same statistical power, and every histogram must be
    filled with the sign of \code{genWeight}.
  \item \textbf{Merged samples are not to be compared by jet counts to
    unmerged ones.} The ``third jet'' of a dijet event in a \PYTHIA
    \code{HardQCD} sample is the shower's; in a merged sample it is a
    matrix-element jet above $Q_{\mathrm{cut}}$. The two agree at low \pt
    and differ at high \pt and wide angle, by construction.
\end{itemize}
\end{pitfalls}

\begin{exercises}
\item The standalone shower can be told to start at the kinematic limit
  instead of the factorisation scale: run
  \code{standalone\_parton\_showering.py --ptmax 396 --repeat 2000} and
  compare the \pt distribution of the hardest branching, and the family
  \pt of \cref{tab:mother_gives_birth:numbers}, with the default run. This
  is the undamped power shower; which of the two is closer to what a
  $2\to3$ matrix element would give at $\pt\approx100\GeV$, and why?
\item From the \code{--history} JSON of the same run, count the events in
  which at least one branching has $\pt>30\GeV$. This fraction is the share
  of events that an FxFx sample with $Q_{\mathrm{cut}}=30\GeV$ would have
  taken from the $\bbbar+1$-parton matrix element rather than from the
  shower.
\item Open a CMS \MG+\PYTHIA fragment for Drell--Yan with FxFx (the
  \code{DYJetsToLL\_M-50\_TuneCP5\_13TeV-amcatnloFXFX-pythia8} dataset) and
  identify \code{JetMatching:qCut}, \code{JetMatching:nJetMax} and
  \code{JetMatching:nQmatch}. Why is \code{nQmatch} five for Drell--Yan and
  what would it have to be for our $\bbbar$ process?
\end{exercises}

\subsection*{Open problems in the context of Phase-2 LHC}
\label{sec:whose_child:open_problems}

The matching and merging of the Run-2 samples was designed for a few
hundred fb$^{-1}$; the $3000\,\mathrm{fb}^{-1}$ of the High-Luminosity
LHC and the multiplicities of its jets strain it in three
ways~\cite{Azzi:2019yne,HSFPhysicsEventGeneratorWG:2020gxw}.
\begin{enumerate}
  \item \textbf{Negative weights and the event budget.} At $20$--$30\%$
    negative weights an NLO-matched sample costs several times the
    simulation of a positive-weight one, and the detector simulation, not
    the generator, is the expensive step. Reweighting the subtraction
    \cite{Frederix:2020trv}, resampling \cite{Andersen:2020sjs} and
    positive-weight matching schemes \cite{Danziger:2021xvr} each trade
    some accuracy for statistical power; which trade is acceptable for a
    Phase-2 $\ttbar$+jets sample is an open question of the HSF generator
    working group.
  \item \textbf{The merging scale as a systematic.} The merging-scale
    dependence of $V$+jets predictions at high jet multiplicity, studied for
    the HL-LHC in \cite{Hoche:2019flt} and across generators in the Les
    Houches comparisons \cite{Amoroso:2020lgh}, is of the same size as the
    experimental precision Phase-2 will reach. Sector showers
    \cite{Brooks:2020mab} and the logarithmically accurate showers of
    \cref{sec:mother_gives_birth:open_problems} promise merged predictions
    whose $Q_{\mathrm{cut}}$ dependence is a genuine higher-order effect
    rather than a shower artefact; none is yet in a CMS production sample.
  \item \textbf{Matching and the heavy-quark mass.} The interpretation of
    the top mass measured from jets depends on how the matched hardest
    emission and the shower treat the quark's mass
    \cite{Nason:2017cxd,Hoang:2020iah}; the same question exists for the
    $b$ quark of our event and its dead cone
    (\cref{sec:mother_gives_birth:dead_cone}). Extending NNLO matching
    \cite{Monni:2019whf} from colour singlets to heavy-quark pairs, and
    deciding what mass the shower should carry, is where a Phase-2
    precision top and $b$-jet programme meets this chapter.
\end{enumerate}
